\chapter{INTRODUÇÃO}
\label{cap:introducao}

\section{Problema e Contexto}

Nos últimos anos, a combinação entre o avanço do poder computacional e o desenvolvimento de algoritmos mais sofisticados tem impulsionado significativamente a área de visão computacional. Essa área da inteligência artificial tem se mostrado essencial em diversas aplicações, desde sistemas de segurança e monitoramento, até diagnósticos médicos e veículos autônomos \cite{szeliski2022computer}. Nesse cenário, o reconhecimento de padrões e a detecção de objetos se tornaram tarefas cada vez mais precisas e eficientes, aproximando as capacidades das máquinas da percepção humana.

O reconhecimento de objetos em imagens é uma atividade simples para humanos quando possuem expertise no assunto e está se tornando cada vez mais acessível aos computadores, devido aos avanços em estudos de visão computacional que utilizam redes neurais de aprendizagem profunda \cite{goodfellow2016deep}. Segundo \citeonline{haykin2009neural}, a aprendizagem profunda é "uma subárea da aprendizagem de máquina que emprega algoritmos inspirados no funcionamento do cérebro humano, simulando um conjunto de neurônios conectados e organizados em camadas, onde cada camada processa informações e transmite o resultado para a próxima".

Com a popularização de arquiteturas como as Redes Neurais Convolucionais (CNNs), tornou-se possível processar grandes volumes de dados visuais com alta acurácia. Essas redes têm a capacidade de extrair automaticamente características relevantes das imagens, dispensando a necessidade de intervenção manual no processo de extração de atributos \cite{lecun2015deep}. Como resultado, surgem soluções mais robustas e adaptáveis, contribuindo diretamente para a expansão de aplicações baseadas em visão computacional.

\section{Objetivos}

O presente trabalho tem como objetivo comparar o desempenho de diferentes arquiteturas de detecção de objetos YOLO, Faster R-CNN, SSD, EfficientDet e RetinaNet utilizando datasets de referência, como COCO e PASCAL VOC. A análise abrangerá métricas quantitativas de desempenho, como precisão, recall, mean Average Precision (mAP) e tempo de inferência, bem como uma discussão crítica sobre a aplicabilidade de cada modelo em diferentes cenários, visando fornecer subsídios para a escolha do algoritmo mais adequado a cada contexto.

\subsection{Objetivo Geral}

O objetivo geral deste trabalho é realizar uma análise comparativa de diferentes arquiteturas de redes neurais artificiais para a tarefa de detecção de objetos, avaliando seu desempenho em termos de acurácia, velocidade e eficiência computacional.

\subsection{Objetivos Específicos}

Para alcançar o objetivo geral proposto, foram definidos os seguintes objetivos específicos:
\begin{itemize}

\item Realizar uma revisão bibliográfica sobre os principais algoritmos de detecção de objetos baseados em redes neurais convolucionais;

\item Selecionar e implementar os modelos YOLO, Faster R-CNN, SSD, EfficientDet e RetinaNet em ambiente TensorFlow;

\item Preparar e pré-processar datasets públicos (COCO e PASCAL VOC) para treinamento e avaliação;

\item Otimizar os hiperparâmetros de cada modelo para garantir uma comparação justa;

\item Treinar os modelos selecionados e avaliar seu desempenho utilizando métricas padronizadas;

\item Coletar e analisar métricas como mAP (mean Average Precision), precisão, recall e FPS (Frames per Second);

\item Realizar uma análise crítica comparativa dos resultados obtidos, considerando as vantagens e limitações de cada arquitetura.

\end{itemize}

\section{Justificativa}

A evolução das arquiteturas de detecção de objetos, como YOLO, Faster R-CNN, SSD, EfficientDet e RetinaNet, impulsionou significativamente a precisão e a velocidade de sistemas de visão computacional \cite{redmon2016yolo, ren2017faster, liu2016ssd, tan2020efficientdet, lin2020focal}. No entanto, apesar dos avanços, cada modelo apresenta características, vantagens e limitações que podem torná-lo mais adequado para determinados contextos e restrições de hardware.

Como afirmam \citeonline{liu2016ssd}, modelos como o SSD se destacam pela alta velocidade de inferência, enquanto arquiteturas como Faster R-CNN \cite{ren2017faster} priorizam a precisão em detrimento do tempo de processamento. Já o YOLO, segundo \citeonline{redmon2016yolo}, "propõe uma abordagem unificada e em tempo real para a detecção de objetos, transformando o problema em uma única tarefa de regressão".

Nesse sentido, torna-se fundamental compreender, de forma comparativa, o desempenho dessas arquiteturas frente a diferentes bases de dados e condições de inferência. Uma análise comparativa criteriosa permite identificar quais modelos oferecem melhor equilíbrio entre precisão e velocidade, considerando não apenas métricas quantitativas, mas também a complexidade de implementação e a escalabilidade. Essa abordagem auxilia na escolha da solução mais adequada e contribui para o avanço do estado da arte.

\section{Estrutura do Trabalho}

Este trabalho está dividido em sete capítulos. O \autoref{cap:introducao} expõe o contexto do estudo, o problema de pesquisa, as justificativas desta pesquisa e os objetivos a serem atingidos. O \autoref{cap:fundamentacao} apresenta a fundamentação teórica sobre redes neurais convolucionais, algoritmos de detecção de objetos e métricas de avaliação. O \autoref{cap:trabalhos} apresenta uma revisão dos trabalhos relacionados na área de detecção de objetos. O \autoref{cap:metodologia} descreve a metodologia empregada, os datasets utilizados e o ambiente experimental. O \autoref{cap:desenvolvimento} detalha a implementação dos modelos e os procedimentos de treinamento. O \autoref{cap:resultados} apresenta e discute os resultados obtidos na análise comparativa. Finalmente, o \autoref{cap:conclusão} apresenta as conclusões e sugere trabalhos futuros.